% ---------------------------------------------------------------------------
%  03-antichain.tex -- the sequential insertion procedure everything preserves
% ---------------------------------------------------------------------------
\section{The antichain update: what one insertion does}
\label{sec:antichain}

Before any parallelism, we need the operation that is being parallelised.  The
passed set is a map
\[
  \dk \;\longmapsto\; A(\dk) \subseteq \text{Zones},
\]
where $A(\dk)$ is a set of pairwise non-subsumed zones -- an \emph{antichain}
-- all belonging to the same discrete state $\dk$.  Inserting a freshly
generated zone $z$ at $\dk$ is \code{NetPassed::add}, and it does five things
(\cref{fig:add}).

\FigAdd

\begin{enumerate}[leftmargin=1.9em, itemsep=3pt]
\item[\textbf{(1)}] \textbf{Compressed fast path.}  For each stored $Z_i$, test
      $[z] \subseteq [Z_i]$ \emph{directly on the packed minimal graph}, in
      $O(|G_i|)$.  On a hit, drop $z$.  By \cref{lem:incl-implies-alu} this is
      a conclusion step~(2) would also have reached, so firing it early changes
      no state count.  It is a pure optimisation of (2) --- but the one that
      matters most, because it decides $89.8\%$ of insertions on
      \code{arbiter17} without ever decompressing a stored zone.

\item[\textbf{(2)}] \textbf{Subsumption, incoming direction.}  If
      $z \subs Z_i$ for some stored $Z_i$, drop $z$.

\item[\textbf{(3)}] \textbf{Subsumption, stored direction.}  Retire every
      stored $Z_i$ with $Z_i \subs z$.

\item[\textbf{(4)}] \textbf{Exact convex-union merge.}  Repeatedly: if some
      live $Z_i$ satisfies $[z \sqcup Z_i] = [z] \cup [Z_i]$ (the convex hull
      adds no new points), set $z := z \sqcup Z_i$ and retire $Z_i$.  Iterate
      until no more merges are possible.

\item[\textbf{(5)}] \textbf{Insert.}  Intern $z$, prepend the new node to the
      chain, stamp it with the current BFS level.
\end{enumerate}

\begin{intuition}
Adding a new shape to the collection stored for one board position.  First:
is the new shape already inside one of the old ones?  Then forget it.  Second:
does the new shape swallow some old ones?  Then delete those.  Third -- and
this is the unusual step -- if the new shape and an old one happen to fit
together into a single shape with no gap, glue them into one.  Finally, store
what is left.
\end{intuition}

\subsection{Why step (4) is the load-bearing one}

Steps (1)--(3) are what every zone-based checker does.  Step (4) is what makes
this design work in parallel, for a reason that is not obvious: it controls the
\emph{length} of the chains, and chain length is exactly the size of the
critical section.

Without step (4), a discrete state accumulates every mutually incomparable zone
it ever sees.  With it, the reachable zone at a discrete state usually
collapses to a single zone: measured mean chain length is $1.11$ on
\code{phil}, $1.00$ on \code{viking}, $1.00$ on \code{arbiter17}.  Since
\code{add} folds over the chain, that is the difference between an $O(1)$ and
an $O(k)$ critical section, and it is what lets \cref{sec:parallel} put a real
zone computation behind a lock without serialising the tool.

It also has a direct effect on state counts, which is where \tool{} beats
\uppaal{} algorithmically rather than by throwing cores at it: on
\code{phil12}, \tool{} stores $802\,871$ symbolic states where \uppaal{} stores
$1\,140\,361$ (\cref{sec:eval}).

The correctness of the test is \cref{thm:hull} in \cref{sec:theory}, and its
geometry is \cref{fig:hull}.  It is worth stressing there that the test is
\emph{exact}: it accepts a merge if and only if the union really is convex.  It
never approximates, so it can never introduce a valuation the system cannot
reach.


One implementation detail of the test matters enough to record, because getting
it wrong is invisible in the results and expensive in the timings.  The test
examines each \emph{strict facet} of either operand, and a facet already implied
by another path need not be examined at all; identifying those is a
minimal-graph analysis costing $O(\dim^3)$.  Running it unconditionally is a
loss whenever there is little to prune, and two zones differing in a handful of
entries then pay the full analysis to discover there was nothing to do.  The
facets are therefore counted first --- an $O(\dim^2)$ scan needing no analysis
--- and the analysis runs only above a threshold
(\code{PTCTLMC\_ANALYZEMIN}, default $8$).  The degenerate case falls out for
free and needs no threshold: if an operand has \emph{no} strict facet then
$Z[i][j] \ge H[i][j]$ everywhere, and since $H$ is the entrywise maximum
$Z[i][j] \le H[i][j]$ always, so $Z = H$ and the hull is that operand.  On
\code{single6} this is worth $-16\%$ with every state count unchanged.

The gate is the facet count and not $\dim$, which is the natural but wrong
choice: \code{csma14} has $\dim = 16$ and is unaffected by the analysis, while
\code{viking20} at $\dim = 12$ is not, so $\dim$ does not predict the cost.

\FigHull

\subsection{Two bounds that make step (4) affordable}

Step (4) restarts its scan after every successful merge, so on a long chain it
is quadratic.  Two caps keep it in hand, and both are sound because
\emph{merging is an optimisation, never an obligation}: skipping a merge only
leaves two zones where one would have done.

\begin{itemize}[leftmargin=1.2em, itemsep=2pt]
\item \textbf{Merge window} (\code{PTCTLMC\_MERGESCAN}, default $32$): only the
      $32$ newest live zones are candidates.  On the pathological \code{single6}
      model (one location, thousands of incomparable zones) the unbounded scan
      ran $4.6$k convex-union tests \emph{per insertion}.

      The obvious justification --- that zones which merge are recent
      neighbours, so a small window loses nothing --- is \emph{false}, and we
      record it because it is the natural thing to assume.  Instrumenting the
      unbounded run on \code{single6} (the only configuration that can see where
      mergeable partners actually are) gives a flat-tailed distribution: a depth
      of $8$ captures $37.5\%$ of merges, $32$ captures $43.6\%$, and $46.4\%$
      lie deeper than $128$.  A window of $8$ is not a cheap approximation of the
      merge set, it misses most of it --- and it is measurably \emph{worse},
      storing $1.65\times$ more zones and running $3.2\times$ slower than $32$.

      What the window really buys is not cheaper merging but shorter chains.
      Fewer merges leave a longer antichain, and every later insertion pays an
      unbounded subsumption scan over it; total scan work (insertions $\times$
      mean chain) is $6.5$M at $32$ against $21.2$M at $8$, a ratio of $3.26$
      that predicts the observed $3.2\times$.  Unbounded is worse again for the
      opposite reason: it buys the shortest chains but pays $28\times$ the
      convex-union tests, and since merging is non-confluent (a hull is a
      different zone with different successors) it explores \emph{more} states,
      not fewer.  $32$ sits where the two costs cross.

\item \textbf{Dense mirror for hot chains} (\code{PTCTLMC\_HOTWORK} /
      \code{PTCTLMC\_HOTMINLEN}, defaults $16$ and $8$): a chain may additionally
      keep its zones in dense canonical form, so the folds read them in place
      instead of decompressing each one.

      The rule is what a chain has \emph{cost}, not how long it is.  A mirror
      saves an $O(L \cdot \dim^3)$ decompression on every visit to a chain of $L$
      live zones and costs $L \cdot \dim^2$ words once, so the payoff is driven
      by $\mathit{visits} \times L$; a length-$16$ chain walked a hundred times
      wastes far more than a length-$64$ chain walked once.  A chain is
      therefore mirrored once $\mathit{visits} \times L \ge 16$, subject to
      $L \ge 8$.

      A fixed length threshold of $64$ --- the earlier rule --- is the
      instructive failure.  It was tuned on \code{single}, whose chains run
      $\sim1000$ zones, and the assumption that "chains that long occur only on
      zone-heavy models" made it look free.  \code{phil} falsifies it: its chains
      sit in the $8$--$63$ range, never reaching the threshold, yet long enough
      that decompressing them on every visit dominates --- $1.14 \times 10^9$
      instructions in \code{mindbm\_unpack}, $11.5\%$ of \code{phil13}.  Moving to
      the work rule is worth $-28\%$ on \code{phil13} and $-27\%$ on \code{phil14}.

      Both halves of the rule are load-bearing.  Hotness alone over-promotes:
      gating on visits only, \code{fischer12} went $0.84 \to 2.48$\,GB, because a
      mirror of a length-$1$ chain is a dense \emph{private} copy of a zone the
      interner was \emph{sharing}, which defeats interning outright.  Length
      alone is the fixed rule that missed \code{phil}.  Promotion is additionally
      bounded by a memory budget (\code{PTCTLMC\_HOTBUDGET}, default a quarter of
      RAM); refusing to promote is always sound, because the mirror is a pure
      cache over packed records that remain the storage of record.
\end{itemize}

\subsection{When exact subsumption is forced}

Two situations make the $\aLU$ test unusable and force plain inclusion:

\begin{itemize}[leftmargin=1.2em, itemsep=2pt]
\item \textbf{Diagonal constraints.}  If the network compares two clocks
      directly ($x - y \prec c$), the $LU$ abstraction is unsound
      \cite{Bouyer04}.  \code{network\_has\_diagonal} detects this once, and the
      engine falls back to exact inclusion plus a diagonal-preserving
      extrapolation (\code{dbm\_extrapolate\_lu\_plus\_nodiag}, which widens only
      the entries touching the reference clock and leaves every true diagonal
      exact).

\item \textbf{Per-valuation properties.}  \code{E<> deadlock} asks about
      individual valuations, not about reaching a location.  $\aLU$ may drop
      valuations, so it is not sound for such a query; the engine sets
      \code{force\_exact} and obtains finiteness from widening instead.  This is
      a genuine parameterisation of one engine, not a second engine: deadlock
      checking used to be a separate sequential search and now shares the
      parallel machinery.
\end{itemize}
